\documentclass{article}

\usepackage{amsmath,amssymb}
\usepackage{xurl}
\usepackage[hidelinks]{hyperref}
\setlength{\emergencystretch}{2em}

% Shared commands for notes in docs/paper-gaps/.

\DeclareUnicodeCharacter{2080}{\ifmmode{}_{0}\else\textsubscript{0}\fi}
\DeclareUnicodeCharacter{2081}{\ifmmode{}_{1}\else\textsubscript{1}\fi}
\DeclareUnicodeCharacter{2082}{\ifmmode{}_{2}\else\textsubscript{2}\fi}
\DeclareUnicodeCharacter{2099}{\ifmmode{}_{n}\else\textsubscript{n}\fi}

\newcommand{\Lean}{\textsc{Lean}}
\ExplSyntaxOn
\NewDocumentCommand{\leanid}{m}{
  \ifmmode
    \textsf{\detokenize{#1}}
  \else
    \group_begin:
    \urlstyle{sf}
    \tl_set:Nn \l_tmpa_tl {#1}
    \tl_replace_all:Nnn \l_tmpa_tl {\_} {_}
    \exp_args:NV \path \l_tmpa_tl
    \group_end:
  \fi
}
\ExplSyntaxOff
\providecommand{\tr}{\operatorname{tr}}
\newcommand{\tnleanIssueBase}{https://github.com/LionSR/TNLean/issues/}
\newcommand{\tnleanPrBase}{https://github.com/LionSR/TNLean/pull/}
\newcommand{\tnleanDocBase}{https://sirui-lu.com/TNLean/docs/}
\newcommand{\ghissue}[1]{\href{\tnleanIssueBase#1}{GitHub issue \##1}}
\newcommand{\ghpr}[1]{\href{\tnleanPrBase#1}{PR \##1}}
\newcommand{\doclink}[1]{\href{\tnleanDocBase#1.html}{\path{#1}}}

% Dirac notation and the identity operator, as in blueprint/src/macros/common.tex.
% \providecommand keeps note-local definitions authoritative.
\usepackage{dsfont}
\providecommand{\ket}[1]{|#1\rangle}
\providecommand{\bra}[1]{\langle #1|}
\providecommand{\braket}[2]{\langle #1 | #2 \rangle}
\providecommand{\ketbra}[2]{|#1\rangle\!\langle #2|}
\providecommand{\Id}{\mathds{1}}
\providecommand{\one}{\mathds{1}}
\providecommand{\C}{\mathbb{C}}
\providecommand{\MN}[1]{M_{#1}(\C)}
\providecommand{\MD}{M_D(\C)}

% External referencing for paper sources.  The build helper writes sanitized
% auxiliary files under build/paper-aux/ and excludes duplicated source labels.
\usepackage{xr}

% Import a generated paper auxiliary file with a caller-supplied label prefix.
% The candidate paths support builds from docs/paper-gaps/, the repository root,
% and build/paper-aux/ itself.
\newcommand{\importpaperaux}[2]{%
  \IfFileExists{../../build/paper-aux/#2.aux}{%
    \externaldocument[#1]{../../build/paper-aux/#2}%
  }{%
    \IfFileExists{build/paper-aux/#2.aux}{%
      \externaldocument[#1]{build/paper-aux/#2}%
    }{%
      \IfFileExists{#2.aux}{%
        \externaldocument[#1]{#2}%
      }{}%
    }%
  }%
}

% General external reference macros
\newcommand{\paperref}[2]{\ref{#1-#2}}
\newcommand{\papereqref}[2]{(\ref{#1-#2})}

% CPSV16 source (1606.00608 / MPDO-22-12-17-2.tex) external document and shortcuts
\importpaperaux{cpsv16-}{1606.00608/MPDO-22-12-17-2-tnlean}
\newcommand{\cpsvref}[1]{\paperref{cpsv16}{#1}}
\newcommand{\cpsveqref}[1]{\papereqref{cpsv16}{#1}}

% Verdict marker: \gapnote{<kind>}{<status>}. Kind is the policy
% classification (clarification, local-correction, scope-restriction,
% unfaithful, false-source, open-gap); status is open, wip (elimination
% underway), resolved, or historical. Machine-read by the site index;
% prints a verdict line.
\newcommand{\gapnote}[2]{\par\noindent\textbf{Verdict:} #1 (#2).\par}


\title{Normalization of the tensor printed in CPSV16 Example 4.12}
\author{}
\date{2026-08-10}

\begin{document}
\maketitle

\section{Printed representative}

Example~4.12 of Cirac--Perez-Garcia--Schuch--Verstraete
(\path{Papers/1606.00608/MPDO-22-12-17-2.tex:932--939}) prints a tensor $M$
with $d=D=2$ and
\[
  1=M_{00}^{00}=M_{00}^{11}=M_{11}^{00}=-M_{11}^{11},
\]
all other components being zero.  In matrix notation this says
\[
  M^{00}=I,
  \qquad
  M^{11}=\sigma_z,
  \qquad
  M^{01}=M^{10}=0,
\]
where $\sigma_z=|0\rangle\langle0|-|1\rangle\langle1|$.  Closing the virtual
indices gives exactly the displayed family
\[
  \rho^{(N)}(M)=I^{\otimes N}+\sigma_z^{\otimes N}
  \qquad (N>0).
\]
Its trace is $2^N$, not $1$.  Thus the printed $M$ is an unnormalized tensor
representative even though the resulting positive operator determines a
normalized density operator after division by $2^N$.

\section{The fixed-representative mismatch}

The physical-trace transfer of the printed tensor is
\[
  \mathcal T_M=M^{00}+M^{11}=I+\sigma_z
    =\begin{pmatrix}2&0\\0&0\end{pmatrix}.
\]
Consequently
\[
  \mathcal T_M^2=2\mathcal T_M
  \qquad\text{but}\qquad
  \mathcal T_M^2\ne\mathcal T_M.
\]
This proves the project's scale-invariant relation with scale $2$, but it does
not prove the literal ZCL equation of Definition~4.2. Indeed, that equation
fails. Moreover, the trace-preserving maps $\mathcal S$ and $\mathcal T$ in
Definition~4.1 force equality of the one-site and two-site traces for every
virtual boundary.  Equivalently, they force literal idempotence
$\mathcal T_M^2=\mathcal T_M$.  The printed representative therefore cannot
satisfy that fixed-tensor renormalization condition.

This does not contradict the normalized density family.  Scaling the local
tensor by $1/2$ gives the separate representative
\[
  \widehat M=\tfrac12 M,
  \qquad
  \mathcal T_{\widehat M}=\tfrac12\mathcal T_M,
  \qquad
  \mathcal T_{\widehat M}^2=\mathcal T_{\widehat M}.
\]
Moreover,
\[
  \rho^{(N)}(\widehat M)
    =2^{-N}\bigl(I^{\otimes N}+\sigma_z^{\otimes N}\bigr),
\]
which is precisely the normalized density family.  This makes
$\widehat M=(1/2)M$ the likely representative intended by the sentence on
line~938 asserting the existence of trace-preserving maps.  The paper does not
print this factor in the tensor components or explicitly change representatives
before that sentence.

\section{Formalization boundary}

The current formalization keeps the two representatives separate.  For the
literal printed $M$ it proves the exact positive-length formula, positivity,
one-site loss of the $\sigma_z$ term, saturation of the area law, source zero
correlation length only in the project's scale-invariant sense, explicit
failure of literal ZCL, and hence failure of the fixed-representative
trace-preserving renormalization condition.

For $\widehat M=(1/2)M$, let $\pi(a,b)=a\oplus b$ and define Kraus operators
\[
  C_{a,b}=|\pi(a,b)\rangle\langle a,b|,
  \qquad
  R_{a,b}=2^{-1/2}|a,b\rangle\langle\pi(a,b)|.
\]
Each parity fibre has two elements, and hence
\[
  \sum_{a,b} C_{a,b}^{\dagger}C_{a,b}=I_4,
  \qquad
  \sum_{a,b} R_{a,b}^{\dagger}R_{a,b}=I_2.
\]
The corresponding completely positive maps are trace preserving.  Since
\[
  \widehat M^{aa}\widehat M^{bb}
    =\tfrac12\widehat M^{\pi(a,b)\pi(a,b)}
\]
and the off-diagonal physical letters vanish, direct substitution proves both
channel equations of Definition~4.1 for every virtual boundary matrix.  These
results are recorded by
\leanid{MPOTensor.CPSVExample412NormalizedRFP.Mhat},
\leanid{MPOTensor.CPSVExample412NormalizedRFP.mpo_Mhat_eq_normalizedMPO_M},
\leanid{MPOTensor.CPSVExample412NormalizedRFP.Mhat_isMPDO}, and
\leanid{MPOTensor.CPSVExample412NormalizedRFP.Mhat_isRFPViaTS}.

This conclusion concerns precisely the two channel equations.  It does not
settle a second normalization tension in the source.  As a doubled-index MPS,
$M$ has two bond-one sectors with physical vectors $(1,0,0,1)$ and
$(1,0,0,-1)$.  Their normal representatives are obtained by division by
$\sqrt2$.  Consequently the two horizontal normal-block weights of
$\widehat M$ are both $1/\sqrt2$.  Neither has modulus one, whereas line~246
assumes that at least one weight does. Thus the proved bare channel predicate
must not be presented as a proof that $\widehat M$ obeys every standing
canonical-form convention in Definition~4.1.  This boundary is also recorded
in
\path{docs/paper-gaps/cpsv16_unit_weight_rfp_scale_tension.tex}.

The independent four-site entropy calculation is also complete.  For the
normalized state, regroup the sites as $A=\{0\}$, $B=\{1,3\}$, and
$C=\{2\}$.  The $AB$ and $BC$ marginals are $I_8/8$, the $B$ marginal is
$I_4/4$, and the full state is uniform on eight even-parity configurations.
Thus
\[
  S(ABC)=3\log 2,
  \qquad S(B)=2\log 2,
  \qquad S(AB)=S(BC)=3\log 2,
\]
so equality in strong subadditivity fails by the strictly positive defect
$\log 2$.  This proves the state-specific Markov obstruction without using the
normalized fixed-point maps.

The full GSNNCH exclusion is already a statement about the normalized family
because \leanid{MPOTensor.IsGSNNCH} quantifies over
\leanid{MPOTensor.normalizedMPO}. It is now proved directly for the printed
tensor. At four sites, the sector sum in Definition~4.8 splits into positive
factors built from the complementary adjacent-bond products
$B_3^{(x)}B_0^{(x)}$ and $B_1^{(x)}B_2^{(x)}$, acting respectively on the
three-site windows $(3,0,1)$ and $(1,2,3)$. Orthogonality of distinct sectors
removes the cross terms in their product, while commutativity within one sector
identifies the product of the two factors with the four-bond sector product.
The resulting positive overlapping factors commute and give the
quantum Markov decomposition
\leanid{MPOTensor.nonempty_quantumMarkovDecomposition_fourCycle_of_isGSNNCHAt},
hence the strong-subadditivity equality
\leanid{MPOTensor.isSSAEquality_fourCycle_of_isGSNNCHAt}. This derivation uses
only Definition~4.8 and the Beigi--Hayden--Jozsa--Petz--Winter structure; it
does not assume zero correlation length, normality, injectivity, primitivity,
a single sector, or completeness of the sector projections.

The identity
\leanid{MPOTensor.CPSVExample412Literal.fourCycleTripartiteState_eq_generic}
shows that this general regrouping is exactly the state used in the literal
entropy calculation. Consequently
\leanid{MPOTensor.CPSVExample412Literal.M_not_isGSNNCH} proves
$\neg\operatorname{IsGSNNCH}(M)$. This discharges the missing and formerly
unlinked four-cycle task tracked in \ghissue{6072} and the missing
Example~4.12 conclusion tracked in \ghissue{6045}. The earlier attempt to pass
through Proposition
\emph{prop4to2} had proof-path drift: that proposition belongs to a different
condition-level implication and is not used here. Transport to the separately
named $\widehat M$ is unnecessary because the GSNNCH predicate already uses
the normalized family.

\section{Verdict}

\textbf{Verdict: false as printed for the literal fixed representative.}
The tensor $M$ printed in Example~4.12 is the formal tensor
\leanid{MPOTensor.CPSVExample412Literal.M}.  Its positive-length operator
formula is proved by
\leanid{MPOTensor.CPSVExample412Literal.rho_eq_finKronecker}, and its
saturation statement by \leanid{MPOTensor.CPSVExample412Literal.M_isSAL}.
The transfer calculation
\leanid{MPOTensor.CPSVExample412Literal.physTraceTransfer_M_sq} has scale
$2$, while
\leanid{MPOTensor.CPSVExample412Literal.physTraceTransfer_M_not_idempotent}
proves failure of literal idempotence.  Therefore
\leanid{MPOTensor.CPSVExample412Literal.M_not_isRFPViaTS} proves that the
fixed-representative RFP assertion is false for the tensor exactly as printed.

The normalized tensor $\widehat M=(1/2)M$ is the likely local correction.  Its
periodic family is identified with the normalized periodic family of $M$ by
tensor-level scalar rescaling and \leanid{MPOTensor.mpo_smul}; it is not an
independent tensor presentation.  The parity Kraus families above prove the
bare trace-preserving channel predicate without additional hypotheses,
discharging the mathematical task tracked in \ghissue{6044}. They do not repair
the incompatible line-246 unit-weight convention, so the formal verdict must
remain separated into the false literal claim, the true normalized channel
equations, the unresolved unit-weight boundary, and the now-proved full GSNNCH
exclusion.

For the normalized four-site state,
\leanid{MPOTensor.CPSVExample412Literal.fourCycleTripartiteState_not_isSSAEquality}
proves the strict strong-subadditivity defect.  Its proof uses the exact
marginals
\leanid{MPOTensor.CPSVExample412Literal.traceC_fourCycleTripartiteState},
\leanid{MPOTensor.CPSVExample412Literal.traceA_fourCycleTripartiteState}, and
\leanid{MPOTensor.CPSVExample412Literal.traceAC_fourCycleTripartiteState},
together with their entropy calculations. This completes \ghissue{6073}.
The general four-cycle Markov implication and the exact coordinate identity
then give
\leanid{MPOTensor.CPSVExample412Literal.M_not_isGSNNCH}, discharging the tasks
tracked in \ghissue{6072} and \ghissue{6045} without hypotheses beyond the
source GSNNCH definition.

\end{document}
